\subsection{Recording Setup}
To capture consistent and analyzable volleyball serves, all recordings were conducted using a fixed baseline camera setup designed to clearly capture the server, ball trajectory, and full court. The camera was positioned at the center of the baseline, approximately 4 m behind the court, mounted on a tripod at a height of 74 in (188 cm). This configuration ensured that both the serve toss and landing zone were visible within the same frame while minimizing parallax distortion. The camera was tilted downward by approximately 5–10°, allowing the near and far court lines to remain visible without excessive floor coverage.

All sessions were recorded at the Rochester Institute of Technology Hale-Andrews Student Life Center volleyball courts under overhead LED lighting. While lighting conditions varied slightly between sessions due to court availability, no recordings required artificial exposure adjustment or post-processing corrections. The only camera setting modification was disabling the phone’s built-in video stabilization to prevent frame warping and maintain consistent spatial alignment across frames.

Recordings were captured at 1920×1080 resolution and 60 frames per second using a Samsung Galaxy S23. The 60 FPS frame rate helped minimize motion blur during high-speed serves, which was crucial for accurate ball detection and trajectory analysis. The tripod distance of 4 m was measured and kept consistent across sessions to ensure uniform perspective and homography mapping later in the pipeline. Each recording session consisted of one or more video clips, each under 15 minutes in duration, which were transferred directly to a desktop computer for processing and annotation.

No external calibration markers were used during recording; instead, court corners and boundary lines visible in the footage were later used for homography-based calibration. This choice simplified setup while still allowing accurate mapping between image coordinates and real-world court positions.

This setup produced a consistent, reproducible viewpoint across all sessions and players, providing a standardized dataset for model training and evaluation in serve detection, trajectory estimation, and landing prediction tasks.

\subsection{Data Organization}
All of the collected data was organized into a standardized directory structure to maintain consistency across sessions, players, and annotation stages. This structure was designed to support both short-term experimentation and long-term reproducibility, ensuring that each stage of the pipeline, from raw recordings to labeled training data, could be traced and reconstructed if needed. The dataset directory (data/) contains all session recordings, extracted frames, annotations, and metadata required for model training and evaluation. At the highest level, the project directory includes five main components: raw and processed videos, extracted frames, annotations, metadata, and training subsets for the models. Each subdirectory follows a uniform naming convention based on session and serve identifiers, allowing for clear correspondence between files and metadata entries.

The raw/ subfolder stores the original session footage as captured by the baseline camera, organized by recording date and session number (e.g., 2025-09-24/session\_1/). Each raw session folder contains one or more .mp4 files directly transferred from the recording device. These raw videos remain unmodified and serve as the master data source for all later stages of processing. Each raw video contains only one player at a time, though individual sessions can include multiple players overall.

Processed serve clips are stored under the processed/ subfolder, which contains serve-level video segments extracted from the raw sessions using the split\_serves.py utility. Each player has a dedicated folder, with sessions nested inside. Serve clips follow a consistent naming convention (serve\_001.mp4, serve\_002.mp4, etc.) to reflect their chronological order within each session. This hierarchy provides a clean separation between players, sessions, and individual serves while allowing straightforward iteration during model inference and annotation. Across all four sessions, a total of 113 serves were extracted and processed into individual clips from six players and four different balls, capturing a range of serving styles while maintaining a consistent visual perspective. 

\emph{These numbers will change later.}

To support frame-level annotation for both ball detection and court segmentation, as well as model training, all serve clips were decomposed into individual frames stored under data/frames/. Frame extraction was performed using the extract\_frames.py script, which automatically matches the frame rate of the source video, preserving the temporal fidelity of the ball trajectory and motion blur characteristics. Across all sessions, a total of over 16,000 frames were extracted. These images form the basis for both bounding-box annotations used in YOLOv8 ball detection training and binary masks used for court segmentation tasks.

The data/annotations/ directory stores all manually and semi-automatically generated labels related to both the ball and the court. Each annotation type corresponds to a distinct subfolder: binary masks for court line segmentation, YOLO-format bounding boxes for ball detection, and JSON files for court corner calibration. Court calibration data are stored under data/annotations/court\_corners/ as one JSON file per session, each containing six manually labeled intersection coordinates corresponding to visible court corners in the recording. Ball annotations were created using CVAT (Computer Vision Annotation Tool)\cite{cvat2023}, which allowed frame-by-frame bounding box labeling with integrated YOLO export. The annotations were stored in CVAT and later moved into the dataset for model training. The exported labels follow the YOLO text format, with one .txt file per frame. This structure integrates directly with the Ultralytics YOLOv8 framework and supports efficient data augmentation during training.

Complementing the visual and annotation data, the data/metadata/ directory provides structured metadata linking sessions, players, and serves. Two CSV files, players.csv and serves.csv, serve as lightweight relational tables for cross-referencing dataset components. players.csv contains one row per player, listing the player’s identifier, name, handedness, and any relevant notes such as skill level or role (e.g., varsity, club, etc.). serves.csv is the primary index for all serve clips in the dataset. Each row records the player name, serve ID, source video path, processed output path, serve start and end frames, and the annotated landing frame. These fields are automatically updated by the serve-splitting and landing-labeling scripts (split\_serves.py and landing\_frame.py), ensuring consistency between raw recordings, processed serve clips, and corresponding annotations.

All data and configuration files are maintained under version control using Git, ensuring that dataset revisions, annotation corrections, and script updates remain tracked over time. This practice is particularly important for reproducibility in experimental pipelines involving iterative model training and evaluation. Large binary files such as videos and frame sequences are tracked through external storage references to prevent repository bloat, while annotation files and metadata remain directly versioned for transparency. The clear directory structure also facilitates automated dataset formatting and validation. Utility scripts such as format\_datasets.py and auto\_label\_and\_upload.py rely on this consistent organization to locate and process files without manual path specification. By enforcing standardized folder names and file formats, the dataset remains scalable and compatible with future extensions, including additional players, camera angles, or labeling types.


\subsection{Annotation Process}
The annotation process established the foundation for all later stages of serve analysis, providing the spatial and temporal ground truth needed for model training, evaluation, and trajectory reconstruction. 

All labeling was conducted in a semi-automated workflow using the open-source CVAT alongside custom Python scripts for data organization and post-processing. CVAT provided an interactive web interface that supported both manual and interpolated bounding box placement, while the scripts handled export conversion and metadata integration.

For ball detection, each serve clip was uploaded to CVAT, where bounding boxes were drawn around every ball on a frame-by-frame basis. CVAT’s interpolation and automatic tracking features were used to propagate bounding boxes across consecutive frames, substantially reducing manual effort while maintaining spatial precision. After annotation, the labels were exported directly in YOLO format, producing one .txt file per frame with normalized center coordinates and bounding-box width and height. These labels were stored alongside the extracted frames within each serve’s directory to match the YOLOv8 training input structure.

After training intermediate YOLOv8 ball detection models, they were used to semi-automate the annotation process. The model was used to output predictions for the clip frames, which were uploaded to CVAT, greatly reducing the amount of time spent manually adjusting bounding boxes. These predictions were then verified before being completed and used for further training.

Serve landing frames were annotated separately using the landing\_frame.py utility. Each serve clip was reviewed manually, and the frame index where the ball first contacted the court was recorded in the serves.csv metadata file under the landing\_frame column. This temporal annotation enables direct evaluation of predicted landing times and supports future trajectory modeling tasks.

Finally, court geometry was annotated once per recording session using the annotate\_court.py script. Each JSON file in data/annotations/court\_corners/ contains four corner coordinates and each side of the center line, used to compute a homography matrix mapping image coordinates to real-world court positions. This calibration step ensures geometric consistency across all serves recorded under the same setup and allows later visualizations, such as landing zone maps, to be expressed in physical units rather than pixels.

Together, these annotation procedures established a unified, verified dataset linking ball position, timing, and court geometry. The resulting labels formed the backbone of the detection, tracking, and trajectory analysis pipeline developed in this project.

\emph{This turned into a lot of writing that I'll definitely reduce later.}